% Owner: TV3 (P5.3b) -- Reviewer: TV1
\subsection{Policy and Value: Hand-Crafted Knowledge versus AlphaGo}
\label{subsec:discussion_policy_value}

The experiments separate two uses of domain knowledge, which correspond to the two networks of AlphaGo~\cite{silver2016mastering}. A \emph{policy} decides which moves are tried: the candidate restriction and pattern ordering of alpha-beta, the softmax sampling of L3, and the prior of progressive bias all play this role. A \emph{value} judges a position without playing it out: the pattern evaluation at the horizon of alpha-beta and the truncated rollout of L4. Our results show that knowledge helps only if it is cheap relative to what it replaces. L3 is the best-informed random policy, yet at equal time it loses to uniform random rollouts, because it runs twelve times fewer simulations. L4 spends the same knowledge on a value estimate after eight moves and is the only policy that wins at equal time. Progressive bias, which puts the knowledge into the selection step instead of the simulations, gives the largest improvement of all UCT variants. AlphaGo makes the same trade on a larger scale: its value network replaces most of the rollout, and its policy network supplies the prior in the selection formula. AlphaGo Zero~\cite{silver2017mastering} drops the rollout altogether.

On Gomoku, alpha-beta with a good evaluation still dominates MCTS (Table~\ref{tab:t7}). The game is decided by short forcing sequences, open threes and fours, which the pattern evaluation recognizes directly and a depth-limited search reads exactly, while random simulations on $225$ cells must sample these sequences by chance. MCTS profits from the same knowledge when it is given it (Tables~\ref{tab:t6} and~\ref{tab:b4}), which suggests that its weakness here comes from the uninformed simulations rather than from the tree search.

\textbf{Limitations.} All Gomoku results use the free-style rule and boards of size $9$ and $15$; the opening is randomized for only two plies. Time-based comparisons depend on the hardware and on how many games run in parallel, so they are compared only within one experiment. Our implementation is in Python, so the absolute numbers of nodes and simulations are far below those of compiled programs, and a faster implementation could shift the balance between cheap and informed rollouts. The never-lose threshold $I^{*}$ is measured on a grid and over $20$ seeds, so it is an empirical bound, not a proof. Finally, the evaluation weights and rollout parameters ($T$, $\tau$, $d_{\mathrm{cut}}$) were set by hand and not tuned.
